\section{Campo, geometría y existencia de soluciones}
\label{sec:ruta-geometrica-libros}

Consideramos tres formas complementarias de estudiar una ecuación dinámica.
La geometría diferencial identifica el espacio de fases y el campo tangente;
la teoría geométrica describe las órbitas, sus límites y sus cuencas; los
métodos topológicos establecen existencia de ceros y puntos fijos mediante
continuidad, compacidad y comportamiento en la frontera. Estas perspectivas
se enlazan cuando pasamos de una ecuación a un retorno o a un operador
integral.

El lenguaje de variedades, formas y Stokes se desarrolla en
\cite{NahmadAchar2022Topologia}. La construcción del flujo, la linealización
y las variedades invariantes corresponden a la teoría geométrica de
\cite{JaurezEtAl2021TGEDI,JaurezEtAl2021TGEDII}. Los métodos de disparo,
puntos fijos y grado de \cite{Amster2021MetodosTopologicos} permiten formular
problemas de existencia aun cuando no dispongamos de una solución explícita.
La relación con los operadores fraccionarios se organiza en
\cref{sec:programa-estudio-bibliografico}.

\subsection{Tres estructuras topológicas de la dinámica}

\begin{definition}[topología del espacio de fases]
Una topología sobre un conjunto \(X\) es una familia \(\mathcal T\) de
subconjuntos, llamados abiertos, que contiene \(\varnothing\) y \(X\), es
cerrada bajo uniones arbitrarias y bajo intersecciones finitas. La pareja
\((X,\mathcal T)\) es un espacio topológico.
\end{definition}

Una métrica induce abiertos mediante bolas. La topología resultante permite
definir continuidad, convergencia y vecindades, pero no conserva una distancia
particular. Si \(h:X\to Y\) es un homeomorfismo, \(h\) y \(h^{-1}\)
transportan abiertos; por eso conservan conexidad y compacidad. Una bola
puede transformarse en un elipsoide sin que cambie esa estructura.

En una variedad diferenciable, el campo \(f(x)\in T_xM\) asigna un vector
tangente a cada punto. Al integrarlo construimos el flujo
\(\varphi_t:M\to M\). La dinámica topológica estudia la acción de esas
aplicaciones sobre conjuntos invariantes y cuencas. Por otra parte, el grado
de un campo o de un residual de retorno estudia la existencia de ceros. El
grado y el flujo se relacionan a través de sus objetos, pero responden a
preguntas distintas.

\begin{center}
\begin{tikzpicture}[
 box/.style={draw=Azul,rounded corners,fill=AzulClaro,align=center,
              minimum width=4.4cm,minimum height=1.25cm},
 arr/.style={-{Latex[length=2.5mm]},thick,draw=GrisOscuro},
 node distance=1.7cm and 1.15cm]
 \node[box] (geom) {Variedades y geometría\\campo $f(x)\in T_xM$};
 \node[box,right=of geom] (flow) {Flujo y dinámica\\$\varphi_t:M\to M$};
 \node[box,right=of flow] (top) {Conjuntos invariantes\\cuencas y retornos};
 \node[box,below=of flow] (exist) {Existencia topológica\\punto fijo y grado};
 \draw[arr] (geom) -- node[above,font=\small]{integración} (flow);
 \draw[arr] (flow) -- node[above,font=\small]{evolución} (top);
 \draw[arr] (exist) -- node[right,font=\small]{existencia de objetos} (flow);
 \draw[arr,dashed] (top) to[bend left=18]
   node[below,font=\small]{residual de retorno} (exist);
\end{tikzpicture}
\end{center}

\subsection{Del cálculo algebraico al objeto geométrico}

Para un equilibrio resolvemos \(f(e)=0\) y calculamos \(Df(e)\).
La matriz describe la evolución de perturbaciones pequeñas, mientras que
la cuenca contiene todos los datos cuyo futuro se aproxima al mismo destino.
Para un ciclo, la ecuación relevante pasa a ser \(P(z)-z=0\) sobre una
sección transversal. Así, cambiar de objeto modifica también el operador
al que se aplica el cálculo algebraico.

\begin{table}[H]
\centering
\caption{Relación entre expresiones algebraicas y objetos dinámicos.}
\label{tab:diccionario-algebra-geometria}
\small
\begin{tabularx}{\textwidth}{@{}L{2.4cm}YYY@{}}
\toprule
Objeto & Expresión & Geometría & Propiedad dinámica \\
\midrule
Campo & $\dot x=f(x)$ & Vector en $T_xM$ & Evolución local del estado. \\
Equilibrio & $f(e)=0$ & Singularidad del campo & Órbita estacionaria. \\
Órbita & $\{\varphi_t(x_0)\}$ & Curva integral orientada & Límites temporales. \\
Linealización & $\dot\xi=Df(e)\xi$ & Dinámica tangente & Estabilidad local. \\
Ciclo & $\varphi_T(x)=x$ & Curva cerrada regular & Periodicidad orbital. \\
Retorno & $P:\Sigma\to\Sigma$ & Cruces sucesivos & Estabilidad transversal. \\
Cuenca & $\mathcal B(A)$ & Datos atraídos por $A$ & Organización de destinos. \\
Ocultedad & $\mathcal B(A)\cap U_e=\varnothing$ & Separación de vecindades &
Inaccesibilidad desde equilibrios próximos. \\
\bottomrule
\end{tabularx}
\end{table}

\subsection{Una silla lineal y sus variedades}

Consideremos
\begin{equation}
 \dot x=x,\qquad \dot y=-y.
 \label{eq:guide-linear-saddle}
\end{equation}
Separar variables da \(\dd x/x=\dd t\) y \(\dd y/y=-\dd t\) para las
componentes no nulas. Al incorporar también las soluciones nulas y el dato
inicial, encontramos
\(\varphi_t(x_0,y_0)=(e^tx_0,e^{-t}y_0)\). Para converger al origen en
tiempo positivo es necesario \(x_0=0\), porque \(|e^tx_0|\to\infty\)
si \(x_0\ne0\); con \(x_0=0\), la segunda componente siempre tiende a
cero. Para tiempo negativo se intercambian los papeles. Por tanto,
\[
 W^s(0)=\{(0,y):y\in\R\},\qquad
 W^u(0)=\{(x,0):x\in\R\}.
\]
La misma fórmula demuestra invariancia de ambos ejes. Fuera de ellos,
\(x(t)y(t)=x_0y_0\), de modo que las órbitas están contenidas en ramas de
hipérbolas. Las variedades estable e inestable organizan así un retrato
calculable, además de proporcionar direcciones tangentes.

\subsection{Una frontera de cuenca explícita}

Para
\begin{equation}
 \dot x=x-x^3=x(1-x^2),
 \label{eq:guide-bistable-line}
\end{equation}
los ceros son \(-1,0,1\), con derivadas \(-2,1,-2\), respectivamente.
Si \(0<x<1\), el campo es positivo; si \(x>1\), es negativo. Por
unicidad las soluciones no cruzan el equilibrio \(1\), por lo que cada
solución positiva es monótona hacia él y permanece acotada. Su límite
\(\ell\) satisface \(\ell-\ell^3=0\): si el campo no se anulara en el
límite, una velocidad de signo fijo impediría la convergencia. El signo del
campo cerca del origen excluye ese punto como destino de un dato positivo.
Por simetría impar, los datos negativos convergen a \(-1\).

También podemos comprobarlo mediante una solución explícita. Para
\(x_0\ne0\), ponemos \(u=x^2\); entonces \(\dot u=2u(1-u)\).
La separación de variables da
\(\log|u/(1-u)|=2t+C\) cuando \(u\ne1\), y conduce a
\[
 x(t)=\frac{\operatorname{sgn}(x_0)}
 {\sqrt{1+(x_0^{-2}-1)e^{-2t}}},\qquad t\ge0.
\]
La expresión incluye los equilibrios \(x_0=\pm1\) por sustitución
directa; \(x_0=0\) da la solución nula. Así,
\(\mathcal B(-1)=(-\infty,0)\), \(\mathcal B(1)=(0,\infty)\) y
\(\{0\}\) es la frontera común. Una malla que no incluya el cero puede
aproximar ambas cuencas sin determinar por sí misma ese punto fronterizo.

\subsection{Fase angular y distancia en el PLL}

Las fases \(\theta\) y \(\theta+2\pi k\), \(k\in\mathbb Z\),
representan el mismo punto. El cociente \(\R/(2\pi\mathbb Z)\) se
identifica con \(S^1\) mediante
\([\theta]\mapsto(\cos\theta,\sin\theta)\). En consecuencia, el espacio
del PLL es \(\R\times S^1\), y la distancia angular es
\begin{equation}
 d_{S^1}(\theta_1,\theta_2)
 =\min_{k\in\mathbb Z}|\theta_1-\theta_2+2\pi k|.
 \label{eq:guide-circle-distance}
\end{equation}
Cambiar un representante angular sólo desplaza el entero que minimiza la
expresión. Para la desigualdad triangular, elijamos enteros minimizadores
\(k_{12}\), \(k_{23}\). La diferencia de la pareja \((1,3)\) con entero
\(k_{12}+k_{23}\) es la suma de las otras dos diferencias; por la
desigualdad triangular real y por la minimización,
\[
 d_{S^1}(\theta_1,\theta_3)\le
 d_{S^1}(\theta_1,\theta_2)+d_{S^1}(\theta_2,\theta_3).
\]
La métrica producto combina esta distancia con la componente real. Un ciclo
corredor avanza \(2\pi\) en el recubrimiento angular y regresa al mismo
punto del cilindro; por eso su aparente apertura en un plano no representa
escape del espacio de fases.

\subsection{Retorno, punto fijo y ciclo}

Sea \(\Sigma\) una sección transversal y \(P\) su primer retorno local.
Si \(P(z)=z\), el tiempo de retorno \(T=\tau(z)>0\) satisface
\(\varphi_T(z)=z\). La propiedad de flujo implica
\(\varphi_{t+T}(z)=\varphi_t(\varphi_T(z))=\varphi_t(z)\); por tanto
la órbita es periódica. La transversalidad excluye un equilibrio en ese
punto. La estabilidad de este ciclo depende de las iteraciones próximas de
\(P\), y por ello de sus multiplicadores transversales, mientras que su
existencia puede establecerse por un teorema de punto fijo o por grado.

En Caputo, el mismo punto de sección puede tener sucesores distintos según
su memoria. El retorno se levanta entonces al espacio de perfiles y la
observación puntual define una relación proyectada. El cambio de espacio
conserva el problema geométrico, pero modifica el dato que determina una
continuación.

\subsection{Conexiones entre las herramientas geométricas}

\begin{longtable}{@{}L{0.8cm}L{3.1cm}L{4.1cm}L{5.55cm}@{}}
\caption{Conexiones entre objetos, fuentes teóricas y sistemas.}
\label{tab:ruta-estudio-geometrica}\\
\toprule
 & Objeto & Fundamento & Aplicación \\
\midrule
\endfirsthead
\toprule
 & Objeto & Fundamento & Aplicación \\
\midrule
\endhead
1 & Continuidad, compacidad y cocientes & Nahmad-Achar & Cilindro del PLL y
distancia angular. \\
2 & Cartas, tangentes y campos & Nahmad-Achar; teoría geométrica I &
Campo como sección de $TM$ y cambios de coordenadas. \\
3 & Existencia, unicidad y flujo & Teoría geométrica I & Regiones suaves y
cruces de Chua PWL. \\
4 & Linealización e hiperbolicidad & Teoría geométrica I--II & Retrato local
de los equilibrios. \\
5 & Variedades invariantes & Teoría geométrica II & Separatrices de Chua,
MAVPD y PLL. \\
6 & Poincaré y Floquet & Teoría geométrica II & Retornos y multiplicadores
de los ciclos PLL y MAVPD. \\
7 & Disparo, punto fijo y grado & Amster & Existencia de ceros del campo y
del residual de retorno. \\
8 & Bifurcaciones y conexiones & Teoría geométrica II & Persistencia de
órbitas y cambios de estabilidad. \\
9 & Dinámica simbólica & Teoría del caos & Herraduras, entropía y crecimiento
de perturbaciones. \\
10 & Cuencas y fronteras & Definiciones de atracción y ocultedad & Separación
respecto de vecindades de equilibrios. \\
11 & Semidinámica de Caputo & Ecuación de Volterra & Perfiles, retornos
funcionales y cuencas de reinicio. \\
\bottomrule
\end{longtable}

\subsection{Problemas sobre los sistemas estudiados}

Para Chua PWL, calcule la solución afín de cada región y relacione el signo
de \(\dot x\) en \(x=\pm1\) con la transversalidad. El campo continuo
produce soluciones únicas, aunque el Jacobiano cambie al atravesar los
planos. En Kalman--Fitts, compare el espectro del equilibrio con la cuenca
del ciclo: el primero controla el entorno local y la segunda incorpora
trayectorias que parten lejos de él. Los 48 sondeos próximos al equilibrio
son una muestra de esa relación, no la descripción del continuo de datos.

Para MAVPD, elija \(H(y)=0\), calcule \(\nabla H\cdot f\) y determine
dónde el retorno es transversal. Relacione la convergencia de los cruces
del ciclo reproducido con la estabilidad de un punto fijo de \(P\). Para
PLL, compare en el recubrimiento y en el cilindro dos fases próximas a
\(-\pi\) y \(\pi\); calcule su distancia angular y explique la
identificación que cierra el ciclo corredor.

\subsection{Aceleración y cambios de coordenadas}

Para \(\dot x=f(x)\) y un difeomorfismo \(h\in C^2\), derivamos dos
veces \(y=h(x)\):
\[
 \dot y=Dh(x)f(x),\qquad
 \ddot y=D^2h(x)[f(x),f(x)]+Dh(x)Df(x)f(x).
\]
Si \(Df(p)f(p)=0\) y \(f(p)\ne0\), el término Hessiano puede ser no
nulo; por tanto, un punto perpetuo clásico depende de la representación
no lineal. En el ejemplo \(\dot x=1\), cualquier punto anula la
aceleración; con \(y=e^x\), tenemos \(\dot y=y\) y \(\ddot y=y>0\).
El cambio es un difeomorfismo de \(\R\) a \((0,\infty)\) y preserva
las órbitas, pero no la aceleración ordinaria. Estudie cómo una conexión
introduce el término adicional necesario para definir aceleración covariante
en una variedad.

\subsection{De la estabilidad local a la cuenca}

La distinción entre trayectoria parametrizada, órbita y flujo aparece ya en
los ejemplos lineales. El grado añade una conclusión de existencia y
Poincaré traduce una órbita periódica en un punto fijo, pero la cuenca se
define mediante el límite de todas las trayectorias que convergen al mismo
conjunto. Para un atractor oculto, esa cuenca evita una vecindad de cada
equilibrio. En la formulación de Caputo, la misma condición depende además
de si comparamos perfiles funcionales o la sección de perfiles constantes
que corresponde a los reinicios convencionales.
